\scanpage{17}

\begin{center}-- 192 --\end{center}
\hfill n$^{\circ}$ 391

2) Le réseau $L_1$ est le sous-$Z$-module de $L_0$ formé des $x=\sum_{i=1}^n\xi_i\varepsilon_i$ tels que
$\sum_{i=1}^n\xi_i$ soit \uline{pair} ; comme $\xi_i$ et $\xi_i^2$ ont même parité, il revient au même de dire
que $(x|x)$ est \uline{pair}. Il est engendré par les $\varepsilon_i\pm\varepsilon_j$, $(i\ne j)$.

3) Le réseau $L_2=L_0+Z.(\frac12\sum_{i=1}^n\varepsilon_i)$ ; il revient au même de dire que les vecteurs
$x=\sum_{i=1}^n\xi_i\varepsilon_i$ de $L_2$ sont ceux qui vérifient les relations :
\[
(4)\qquad 2\xi_i\in Z,\quad \xi_i-\xi_j\in Z\quad\text{quels que soient }i,j.
\]
Comme $(\varepsilon_k|\frac12\sum_i\varepsilon_i)=\frac12$ pour tout $k$, et
$\|\frac12\sum_i\varepsilon_i\|^2=n/4$, on a $(u|v)\in\frac12 Z$
pour $u,v$ dans $L_2$ si $n$ est pair.

4) Si $n$ est \uline{multiple de} 4, le réseau $L_3=L_1+Z.(\frac12\sum_{i=1}^n\unsure{\xi_i|\varepsilon_i})$, sous-$Z$-module de $L_2$;
c'est l'ensemble des vecteurs $x=\sum_{i=1}^n\xi_i\varepsilon_i$ de $V$ vérifiant les relations (4) et
en outre tels que $\sum_{i=1}^n\xi_i\in 2Z$ ; on a $(u|v)\in Z$ quels que soient $u,v$ dans $L_3$.

On notera que le réseau \uline{associé} à $L_0(V)$ est $L_0(V^*)$, le réseau \uline{associé}
à $L_1(V)$, (resp. $L_2(V)$), est $L_2(V^*)$, (resp. $L_1(V^*)$) ; enfin, le réseau \uline{associé}
à $L_3(V)$ est $L_3(V^*)$, car c'est le sous-réseau de $L_2(V^*)$ formé des $x=\sum_{i=1}^n\xi_i\varepsilon_i^*$
tels que $\frac12(\sum_{i=1}^n\xi_i)\in Z$, et on vient de voir, (en vertu de l'hypothèse $n\equiv0\ (\mathrm{mod}.4)$)
que ces conditions caractérisaient $L_3(V^*)$.

Nous allons voir successivement, dans chacun des 8 n$^{os}$ suivants, comment
cette méthode démontre en premier lieu l'existence des systèmes de racines
correspondant aux 8 types de schémas de Dynkin autres que $G_2$ ; en outre, pour
chaque type de schéma, nous expliciterons :

\mspage{16}
